\documentclass[11pt]{article}

% ---------- Packages ----------
\usepackage[margin=1in]{geometry}
\usepackage{amsmath, amssymb}
\usepackage{graphicx}
\usepackage{fancyhdr}

% ---------- Paragraph Formatting ----------
\setlength{\parindent}{0pt}
\setlength{\parskip}{0.5em}

% ---------- Header ----------
\pagestyle{fancy}
\fancyhf{}
\lhead{ASEN 5335 Aerospace Environment}
\rhead{Justin Le}
\cfoot{\thepage}

% ---------- Document ----------
\begin{document}

% ---------- Title ----------
\begin{center}
    {\Large \textbf{Homework 1}}\\[0.5em]
    {\large ASEN 5335 Aerospace Environment}\\[0.5em]
    Justin Le\\
    \today
\end{center}

\vspace{1em}

% ============================================================
% Problem 1
% ============================================================

\section*{Problem 1: Radiation from Corona}

\subsection*{(i)}

Given $T$ of the corona region to be $T=2MK$. We can use the equation,
\[
\lambda_{max}T = \alpha
\]

Where $\alpha = \text{constant}(2.98*10^-3 K)$.

Rearrange to solve for wavelength.
\[
\lambda_{max} = \frac{T}{\alpha}
\]

The corona region's temperature can be substituted and solved.
\[
\boxed{\lambda_{max}  = 1.449*10^-9 \, \mathrm{m} = 14.49*10^-10 \, \mathrm{m}}
\]

This would also be in the same range for X-rays, $\approx 10^-10 \, \mathrm{m}$.

\subsection*{(ii)}
For the region photosphere region at a cooler $T=5,778 \, \mathrm{K}$,
the process is the same.

\[
\boxed{\lambda_{max}  = 5.01*10^-7 \, \mathrm{m} = 50.15*10^-8 \, \mathrm{m}}
\]

This would be in the same range for UV $\approx 10^-8 \, \mathrm{m}$.

% ============================================================
% Problem 2
% ============================================================

\section*{Problem 2: Thermal Escape}

From Wikipedia we can obtain an equation for most probable thermal velocity
of a particle.

\[
\begin{aligned}
v_{\mathrm{th}} &= \sqrt{\frac{2k_B T}{m}} \\[4pt]
k_B &= 1.38 \times 10^{-23}\,\mathrm{J/K} \\[4pt]
m &= m_{\mathrm{proton}} = 1.673 \times 10^{-27}\,\mathrm{kg} \\[4pt]
T &= 2\,\mathrm{MK} = 2 \times 10^6\,\mathrm{K} \\
&\boxed{\therefore v_{th} = 181,644 \, \mathrm{m/s}}
\end{aligned}
\]

We then calculate the escape velocity using the equation,
\[
\begin{aligned}
v_{escape} = \sqrt{ \frac{2GM}{r} } \\
G = 6.674*10^{-11} \, \mathrm{Nm^2/kg^2} \\
M = 1.989*10^{30} \, \mathrm{kg} \\
r = 10*10^{6}*10^3 \, \mathrm{m}
\end{aligned}
\]

Where $G$ is the unversial gravitational constant, $M$ is
the mass of the sun, and $r$ is the radius of the of the outer
corona. We then find,

\[
\boxed{v_{escape} = 163,000 \, \mathrm{m/s}}
\]

\[
\begin{gathered}
\because v_{th} < v_{escape} \\
\boxed{\text{It is possible for the coronal protons to escape the sun's gravity}}
\end{gathered}
\]

% ============================================================
% Problem 3
% ============================================================

\section*{Problem 3}

\subsection*{(a)}

Write your solution here.

\subsection*{(b)}

Write your solution here.

% ============================================================
% Problem 4
% ============================================================

\section*{Problem 4}

Write your solution here.

% ============================================================

\end{document}
